% Options for packages loaded elsewhere
\PassOptionsToPackage{unicode}{hyperref}
\PassOptionsToPackage{hyphens}{url}
\PassOptionsToPackage{dvipsnames,svgnames,x11names}{xcolor}
%
\documentclass[
  twocolumn,
  landscape]{report}

\usepackage{amsmath,amssymb}
\usepackage{lmodern}
\usepackage{iftex}
\ifPDFTeX
  \usepackage[T1]{fontenc}
  \usepackage[utf8]{inputenc}
  \usepackage{textcomp} % provide euro and other symbols
\else % if luatex or xetex
  \usepackage{unicode-math}
  \defaultfontfeatures{Scale=MatchLowercase}
  \defaultfontfeatures[\rmfamily]{Ligatures=TeX,Scale=1}
\fi
% Use upquote if available, for straight quotes in verbatim environments
\IfFileExists{upquote.sty}{\usepackage{upquote}}{}
\IfFileExists{microtype.sty}{% use microtype if available
  \usepackage[]{microtype}
  \UseMicrotypeSet[protrusion]{basicmath} % disable protrusion for tt fonts
}{}
\makeatletter
\@ifundefined{KOMAClassName}{% if non-KOMA class
  \IfFileExists{parskip.sty}{%
    \usepackage{parskip}
  }{% else
    \setlength{\parindent}{0pt}
    \setlength{\parskip}{6pt plus 2pt minus 1pt}}
}{% if KOMA class
  \KOMAoptions{parskip=half}}
\makeatother
\usepackage{xcolor}
\setlength{\emergencystretch}{3em} % prevent overfull lines
\setcounter{secnumdepth}{-\maxdimen} % remove section numbering
% Make \paragraph and \subparagraph free-standing
\ifx\paragraph\undefined\else
  \let\oldparagraph\paragraph
  \renewcommand{\paragraph}[1]{\oldparagraph{#1}\mbox{}}
\fi
\ifx\subparagraph\undefined\else
  \let\oldsubparagraph\subparagraph
  \renewcommand{\subparagraph}[1]{\oldsubparagraph{#1}\mbox{}}
\fi


\providecommand{\tightlist}{%
  \setlength{\itemsep}{0pt}\setlength{\parskip}{0pt}}\usepackage{longtable,booktabs,array}
\usepackage{calc} % for calculating minipage widths
% Correct order of tables after \paragraph or \subparagraph
\usepackage{etoolbox}
\makeatletter
\patchcmd\longtable{\par}{\if@noskipsec\mbox{}\fi\par}{}{}
\makeatother
% Allow footnotes in longtable head/foot
\IfFileExists{footnotehyper.sty}{\usepackage{footnotehyper}}{\usepackage{footnote}}
\makesavenoteenv{longtable}
\usepackage{graphicx}
\makeatletter
\def\maxwidth{\ifdim\Gin@nat@width>\linewidth\linewidth\else\Gin@nat@width\fi}
\def\maxheight{\ifdim\Gin@nat@height>\textheight\textheight\else\Gin@nat@height\fi}
\makeatother
% Scale images if necessary, so that they will not overflow the page
% margins by default, and it is still possible to overwrite the defaults
% using explicit options in \includegraphics[width, height, ...]{}
\setkeys{Gin}{width=\maxwidth,height=\maxheight,keepaspectratio}
% Set default figure placement to htbp
\makeatletter
\def\fps@figure{htbp}
\makeatother

\makeatletter
\makeatother
\makeatletter
\makeatother
\makeatletter
\@ifpackageloaded{caption}{}{\usepackage{caption}}
\AtBeginDocument{%
\ifdefined\contentsname
  \renewcommand*\contentsname{Table of contents}
\else
  \newcommand\contentsname{Table of contents}
\fi
\ifdefined\listfigurename
  \renewcommand*\listfigurename{List of Figures}
\else
  \newcommand\listfigurename{List of Figures}
\fi
\ifdefined\listtablename
  \renewcommand*\listtablename{List of Tables}
\else
  \newcommand\listtablename{List of Tables}
\fi
\ifdefined\figurename
  \renewcommand*\figurename{Figure}
\else
  \newcommand\figurename{Figure}
\fi
\ifdefined\tablename
  \renewcommand*\tablename{Table}
\else
  \newcommand\tablename{Table}
\fi
}
\@ifpackageloaded{float}{}{\usepackage{float}}
\floatstyle{ruled}
\@ifundefined{c@chapter}{\newfloat{codelisting}{h}{lop}}{\newfloat{codelisting}{h}{lop}[chapter]}
\floatname{codelisting}{Listing}
\newcommand*\listoflistings{\listof{codelisting}{List of Listings}}
\makeatother
\makeatletter
\@ifpackageloaded{caption}{}{\usepackage{caption}}
\@ifpackageloaded{subcaption}{}{\usepackage{subcaption}}
\makeatother
\makeatletter
\@ifpackageloaded{tcolorbox}{}{\usepackage[many]{tcolorbox}}
\makeatother
\makeatletter
\@ifundefined{shadecolor}{\definecolor{shadecolor}{rgb}{.97, .97, .97}}
\makeatother
\makeatletter
\makeatother
\ifLuaTeX
  \usepackage{selnolig}  % disable illegal ligatures
\fi
\IfFileExists{bookmark.sty}{\usepackage{bookmark}}{\usepackage{hyperref}}
\IfFileExists{xurl.sty}{\usepackage{xurl}}{} % add URL line breaks if available
\urlstyle{same} % disable monospaced font for URLs
\hypersetup{
  pdftitle={Physics Validation of Y3 Clusters Code},
  pdfauthor={Y3 Cluster Team (3ct)},
  colorlinks=true,
  linkcolor={blue},
  filecolor={Maroon},
  citecolor={Blue},
  urlcolor={Blue},
  pdfcreator={LaTeX via pandoc}}

\title{Physics Validation of Y3 Clusters Code}
\author{Y3 Cluster Team (3ct)}
\date{}

\begin{document}
\maketitle
\ifdefined\Shaded\renewenvironment{Shaded}{\begin{tcolorbox}[borderline west={3pt}{0pt}{shadecolor}, enhanced, sharp corners, interior hidden, frame hidden, breakable, boxrule=0pt]}{\end{tcolorbox}}\fi

\renewcommand*\contentsname{Table of contents}
{
\hypersetup{linkcolor=}
\setcounter{tocdepth}{0}
\tableofcontents
}
\hypertarget{introduction}{%
\chapter{Introduction}\label{introduction}}

This document presents demonstrations of the validity of calculations
done in the Y3 clusters analysis code. Each section contains the
validation testing results for a single function in our calculation.

\begin{enumerate}
\def\labelenumi{\alph{enumi}.}
\tightlist
\item
  For each of these functions, we have evaluations from a trusted source
  that are assumed to be correct. We test \emph{our implementation} by
  comparing the results to the evaluations from the trusted source.
\item
  Unless otherwise indicated, we require fractional differences to be
  less than \(\ensuremath{10^{-6}}\), \textbf{or} absolute differences
  less than \(\ensuremath{10^{-12}}\). The relative error bound was
  chosen to be smaller than the systematic uncertainty to which the
  value is known from the trusted source. An absolute error bound is
  needed for the cases in which the predicted value is 0 or very close
  to 0.
\item
  In each section, we show both absolute and residual comparisons. In
  the absolute comparison plots, the \emph{expected} (or \emph{true})
  values are shown in black, and results from our \emph{implementation}
  (or \emph{test} values) are shown in cyan. We also show the fractional
  (relative) and absolute residuals.
\item
  In the plot of the fractional residuals, the color of each point is
  determined by whether or not that point passes the absolute residual
  test. In the plot of the absolute residuals, the color of each point
  is determined by whether or not that point passes the fractional
  residual test. Thus, in each plot, a point passes the test if
  \emph{either} it is blue or within the error bounds (indicated by
  horizontal lines). Points that fail the test will be both red and
  outside of the allowable error bounds (indicated by horizontal lines).
\item
  For a test to pass, it is required that every calculated point pass.
  This condition is verified automatically and the result is noted in
  bold in each section.
\end{enumerate}

\hypertarget{survey-area}{%
\chapter{Survey Area}\label{survey-area}}

\hypertarget{sdss-survey-footprint-function}{%
\section{SDSS Survey footprint
function}\label{sdss-survey-footprint-function}}

\(\Omega_z\) is the area of the survey as a function of redshift \(z\).
This section tests our implementation of the SDSS survey area function,
which is implemented in the file \texttt{omega\_z\_sdss.hh}.

For this testing, we compare our function output to the independent
function evaluations from Matteo Costanzi in the SDSS cluster cosmology
analysis (Costanzi et al.~2019).

\begin{enumerate}
\def\labelenumi{\alph{enumi}.}
\tightlist
\item
  The comparisons are done at 550 unique \(z\) points.
\item
  Over the range tested \((0.05 < z < 0.598)\), the maximum magnitude of
  the absolute residual between the expected value and our approximation
  is \(\ensuremath{6.0902368\times 10^{-7}}\); the maximum magnitude of
  the fractional residual is \(\ensuremath{1.9964541\times 10^{-7}}\).
\item
  For this test, we require the relative residuals to be less than
  \(\ensuremath{10^{-6}}\), or the absolute residuals to be less than
  \(\ensuremath{10^{-12}}\).
\item
  \textbf{For this test, 550 of 550 points pass.}
\end{enumerate}

\begin{figure}

{\centering \includegraphics{physics_validation_files/figure-pdf/unnamed-chunk-3-1.pdf}

}

\caption{Absolute comparison of our implementation results (cyan points)
to the expected values (black points) for \(\Omega_z\). Note that
over-plotting may obscure the expected values.}

\end{figure}

\begin{figure}

{\centering \includegraphics{physics_validation_files/figure-pdf/unnamed-chunk-4-1.pdf}

}

\caption{Relative residual between our implementation results and the
expected values. Blue horizontal lines are the error tolerance bounds.
Blue points pass, and red points fail, the absolute residual test.}

\end{figure}

\hypertarget{halo-mass-function}{%
\chapter{Halo Mass Function}\label{halo-mass-function}}

\(n(M_{200m}, z)\) is the halo mass function, which describes the 3D
cosmic density of halos at a given mass \(M_{200m}\) and \(z\). We have
modified the CosmoSIS halo mass function module to consider the effect
of massive neutrinos (Costanzi et al.~2013)- we use the \(P_{cdm}\)
power spectrum rather than the \(P_{tot}\).

For this testing, we compare results from our implementation to the
independent function evaluations from Matteo Costanzi in the SDSS
cluster cosmology analysis (Costanzi et al.~2019)

\begin{enumerate}
\def\labelenumi{\alph{enumi}.}
\tightlist
\item
  The comparisons are made at 2 unique \(z\) values with 2 different
  mass values.
\item
  Over the range tested (\(12 < M < 15.96\)), the maximum magnitude of
  the absolute residual between the expected value and our calculated
  value is \(\ensuremath{6.128\times 10^{-6}}\); the maximum magnitude
  of the fractional residual is \(0.0055699\).
\item
  For this test, we require the relative differences to be less than
  \(0.01\), or the absolute difference to be less than
  \(\ensuremath{10^{-16}}\).
\item
  \textbf{For this test, 200 of 200 points pass.}
\end{enumerate}

Each panel in the figure below corresponds to a different \(z\) value,
specified in the panel top label.

\begin{figure}

{\centering \includegraphics{physics_validation_files/figure-pdf/unnamed-chunk-7-1.pdf}

}

\caption{Absolute comparison of our implementation to the expected
results. Note that over-plotting may obscure the expected values, if the
approximated values are very close.}

\end{figure}

\begin{figure}

{\centering \includegraphics{physics_validation_files/figure-pdf/unnamed-chunk-8-1.pdf}

}

\caption{Relative residual comparison of our implementation to the
expected results. Blue points pass, and red points fail, the absolute
residual test.}

\end{figure}

\textbf{This test requires running masses\_test.ini in cosmosis\_tests
first} Need to do better than this in the future.

\hypertarget{volume-element}{%
\chapter{Volume Element}\label{volume-element}}

\(\frac{dV}{d\Omega dz}\) is the co-moving volume element in solid angle
\(d \Omega\) and redshift interval \(dz\), and is a function of \(z\).
Our implementation is in the file \texttt{dv\_do\_dz\_t.hh}.

For this testing, we compare our function output to the independent
function evaluations from Matteo Costanzi in the SDSS cluster cosmology
analysis (Costanzi et al.~2019), calculated with \(\Omega_m =0.3\),
\(\Omega_\Lambda =0.7\), \(h_0=0.7\).

\begin{enumerate}
\def\labelenumi{\alph{enumi}.}
\tightlist
\item
  The comparisons are done at 100 unique \(z\) values.
\item
  Over the range tested \((0.01 < z < 0.99)\), the maximum magnitude of
  the absolute residual between the expected value and our approximation
  is \(\ensuremath{2.150846\times 10^{6}}\); the maximum magnitude of
  the fractional residual is \(\ensuremath{2.3831532\times 10^{-4}}\).
\item
  For this test, we require the relative residuals to be less than
  \(0.001\), or the absolute residuals to be less than
  \(\ensuremath{10^{-12}}\).
\item
  \textbf{For this test, 100 of 100 points pass.}
\end{enumerate}

\begin{figure}

{\centering \includegraphics{physics_validation_files/figure-pdf/unnamed-chunk-11-1.pdf}

}

\caption{Absolute comparison of our approximation to the expected
results. Note that over-plotting may obscure the expected values, if the
approximated values are very close.}

\end{figure}

\begin{figure}

{\centering \includegraphics{physics_validation_files/figure-pdf/unnamed-chunk-12-1.pdf}

}

\caption{Relative residual comparison or implementation to the expected
results. Points in blue pass, and points in red fail, the absolute
residual test.}

\end{figure}

\hypertarget{scaling-relations}{%
\chapter{Scaling Relations}\label{scaling-relations}}

\hypertarget{mass-observable-relation}{%
\section{Mass-observable relation}\label{mass-observable-relation}}

\(P(\lambda |M_{200m})\) describes the probabilistic distribution of the
cluster observable richness (\(\lambda\)) at a given mass \(M_{200m}\).
Our calculation is implemented in the file \texttt{mor\_t2.hh}.

For this testing, we compare results from our implementation to the
independent function evaluations from Matteo Costanzi in the SDSS
cluster cosmology analysis (Costanzi et al.~2019)

\begin{enumerate}
\def\labelenumi{\alph{enumi}.}
\tightlist
\item
  The comparisons are made at 51 unique \(\lambda\) values with 6
  different mass values.
\item
  Over the range tested (\(1 < \lambda < 101\)), the maximum magnitude
  of the absolute residual between the expected value and our calculated
  value is \(\ensuremath{1.8422198\times 10^{-8}}\); the maximum
  magnitude of the fractional residual is
  \(\ensuremath{1.995633\times 10^{-5}}\).
\item
  For this test, we require the relative differences to be less than
  \(\ensuremath{5\times 10^{-6}}\), or the absolute difference to be
  less than \(\ensuremath{10^{-12}}\).
\item
  \textbf{For this test, 306 of 306 points pass.}
\end{enumerate}

Each panel in the figure below corresponds to a different
\(\mathrm{log} M_{200m}\) value, specified in the panel top label.

\begin{figure}

{\centering \includegraphics{physics_validation_files/figure-pdf/unnamed-chunk-15-1.pdf}

}

\caption{Absolute comparison of our implementation to the expected
results. Note that over-plotting may obscure the expected values, if the
approximated values are very close.}

\end{figure}

\begin{figure}

{\centering \includegraphics{physics_validation_files/figure-pdf/unnamed-chunk-16-1.pdf}

}

\caption{Relative residual comparison of our implementation to the
expected results. Blue points pass, and red points fail, the absolute
residual test.}

\end{figure}

\begin{figure}

{\centering \includegraphics{physics_validation_files/figure-pdf/unnamed-chunk-17-1.pdf}

}

\caption{Absolute residual comparison of our approximation to the
expected results for MOR. Blue points pass, and red points fail, the
relative residual test.}

\end{figure}

\hypertarget{projection-effect-on-scaling-relations}{%
\section{Projection effect on Scaling
relations}\label{projection-effect-on-scaling-relations}}

Projections of LOS structures cause a shift and an extra scatter to the
values of cluster richness \(\lambda\). \(P(\lambda_c |\lambda_t)\)
describes the probabilistic distribution of the richness scatter
(between \(\lambda_c\) and \(\lambda_t\)) caused by projections. Our
calculation is implemented in the file \texttt{lc\_lt\_t.hh}.

For this testing, we compare results from our implementation to the
independent function evaluations from Matteo Costanzi's projection
analyses (Costanzi et al.~2019)

\begin{enumerate}
\def\labelenumi{\alph{enumi}.}
\tightlist
\item
  The comparisons are made at 10 unique \(\lambda_{c}\) values with 10
  different \(\lambda_{t}\) values.
\item
  Over the range tested (\(2 < \lambda_{cen} < 150\)), the maximum
  magnitude of the absolute residual between the expected value and our
  calculated value is \(\ensuremath{4.8340653\times 10^{-10}}\); the
  maximum magnitude of the fractional residual is
  \(\ensuremath{1.466778\times 10^{-6}}\).
\item
  For this test, we require the relative differences to be less than
  \(\ensuremath{10^{-6}}\), or the absolute difference to be less than
  \(\ensuremath{10^{-12}}\).
\item
  \textbf{For this test, 1001 of 1001 points pass.}
\end{enumerate}

Each panel in the figure below corresponds to a \(\lambda_{t}\) value,
specified in the panel top label.

\begin{figure}

{\centering \includegraphics{physics_validation_files/figure-pdf/unnamed-chunk-20-1.pdf}

}

\caption{Absolute comparison of our implementation to the expected
results. Note that over-plotting may obscure the expected values, if the
approximated values are very close.}

\end{figure}

\begin{figure}

{\centering \includegraphics{physics_validation_files/figure-pdf/unnamed-chunk-21-1.pdf}

}

\caption{Relative residual comparison of our implementation to the
expected results. Blue points pass, and red points fail, the absolute
residual test.}

\end{figure}

\begin{figure}

{\centering \includegraphics{physics_validation_files/figure-pdf/unnamed-chunk-22-1.pdf}

}

\caption{Absolute residual comparison of our implementation to the
expected results. Blue points pass, and red points fail, the relative
residual test.}

\end{figure}

\hypertarget{mis-centering-effect-on-scaling-relations}{%
\section{Mis-centering effect on scaling
relations}\label{mis-centering-effect-on-scaling-relations}}

Mis-centering cause a shift and an extra scatter to the values of
cluster richness \(\lambda\). \(P(\lambda_o/\lambda_c | R_{mis})\)
describes the probabilistic distribution of the richness shift
(\(\lambda_o/\lambda_c\)) with a mis-centering distance \(R_{mis}\) Our
calculation is implemented in the file \texttt{lo\_lc\_t.hh}.

For this testing, we compare results from our implementation to the
independent function evaluations from Yuanyuan Zhang's mis-centering
analysis (Zhang et al.~2019).

\begin{enumerate}
\def\labelenumi{\alph{enumi}.}
\tightlist
\item
  The comparisons are made at 2092 unique values of
  \(\lambda_{o}/\lambda_{c}\) with 20 different \(R_{mis}\) values.
\item
  Over the range tested (\(0.1 < \lambda_{cen} < 2.6457143\)), the
  maximum magnitude of the absolute residual between the expected value
  and our calculated value is \(\ensuremath{8.8817842\times 10^{-15}}\);
  the maximum magnitude of the fractional residual is
  \(\ensuremath{9.2157402\times 10^{-14}}\).
\item
  For this test, we require the relative differences to be less than
  \(\ensuremath{10^{-6}}\), or the absolute difference to be less than
  \(\ensuremath{10^{-12}}\).
\item
  \textbf{For this test, 38542 of 38542 points pass.}
\end{enumerate}

Each panel in the figure below corresponds to a \(R_{mis}\) value,
specified in the panel top label.

\begin{figure}

{\centering \includegraphics{physics_validation_files/figure-pdf/unnamed-chunk-25-1.pdf}

}

\caption{Absolute comparison of our implementation to the expected
results. Note that over-plotting may obscure the expected values, if the
approximated values are very close.}

\end{figure}

\begin{figure}

{\centering \includegraphics{physics_validation_files/figure-pdf/unnamed-chunk-26-1.pdf}

}

\caption{Relative residual comparison of our implementation to the
expected results. Blue points pass, and red points fail, the absolute
residual test.}

\end{figure}

\begin{figure}

{\centering \includegraphics{physics_validation_files/figure-pdf/unnamed-chunk-27-1.pdf}

}

\caption{Absolute residual comparison of our implementation to the
expected results. Blue points pass, and red points fail, the relative
residual test.}

\end{figure}

\hypertarget{number-counts-and-masses-in-sdss}{%
\chapter{Number Counts and Masses in
SDSS}\label{number-counts-and-masses-in-sdss}}

\textbf{This test requires running sdss\_analysis.ini in cosmosis\_tests
first}.

In this section, we compare the stacked number counts and mass
predictions in stacked richness/mass bins for a SDSS like analysis.
These values are the integrated results using components listed above.
For this testing, we compare results from our implementation to the
implementations from Matteo Costanzi for the SDSS cosmology analysis
(Costanzi et al.~2019). The test is conducted for the following set of
cosmology and scaling relation parameters specified in
\textbf{cosmosis\_tests/values\_comp.ini}.

\begin{enumerate}
\def\labelenumi{\alph{enumi}.}
\tightlist
\item
  The comparisons are made at 5 unique values of \(\lambda\) bins.
\item
  Over the \textbf{number of bins} tested and for the stacked
  \textbf{number counts}, the maximum magnitude of the absolute residual
  between the expected value and our calculated value is \(23.515351\);
  the maximum magnitude of the fractional residual is \(0.0229132\).
\item
  Over the \textbf{number of bins} tested and for the stacked
  \textbf{average masses}, the maximum magnitude of the absolute
  residual between the expected value and our calculated value is
  \(0.010556\); the maximum magnitude of the fractional residual is
  \(\ensuremath{7.4472272\times 10^{-4}}\).
\end{enumerate}

\begin{figure}

{\centering \includegraphics{physics_validation_files/figure-pdf/unnamed-chunk-30-1.pdf}

}

\caption{Absolute comparison of our implementation to the expected
results. Note that over-plotting may obscure the expected values, if the
approximated values are very close.}

\end{figure}

\begin{figure}

{\centering \includegraphics{physics_validation_files/figure-pdf/unnamed-chunk-31-1.pdf}

}

\caption{Relative residual implementation or our approximation to the
expected results. Points in blue pass, and points in red fail, the
absolute residual test.}

\end{figure}

\begin{figure}

{\centering \includegraphics{physics_validation_files/figure-pdf/unnamed-chunk-32-1.pdf}

}

\caption{Absolute comparison of our implementation to the expected
results for average masses. Note that over-plotting may obscure the
expected values, if the approximated values are very close.}

\end{figure}

\begin{figure}

{\centering \includegraphics{physics_validation_files/figure-pdf/unnamed-chunk-33-1.pdf}

}

\caption{Relative residual implementation or our approximation to the
expected results for average masses. Points in blue pass, and points in
red fail, the absolute residual test.}

\end{figure}

\hypertarget{number-counts-and-masses-in-desy1}{%
\chapter{Number Counts and Masses in
DESY1}\label{number-counts-and-masses-in-desy1}}

\textbf{This test requires running y1\_analysis.ini in cosmosis\_tests
first}

In this section, we compare the stacked number counts and mass
predictions in stacked richness/mass bins for a y1 like analysis. These
values are the integrated results using components listed above. For
this testing, we compare results from our implementation to the
implementations from Matteo Costanzi for the y1 cosmology analysis (DES
2019). The test is conducted for the following set of cosmology and
scaling relation parameters specified in
\textbf{cosmosis\_tests/values\_comp\_y1.in}

\begin{enumerate}
\def\labelenumi{\alph{enumi}.}
\tightlist
\item
  The comparisons are made at 4 unique values of \(\lambda\) bins and 3
  unique values of \(z\) bins.
\item
  Over the \textbf{number of bins} tested and for the stacked
  \textbf{number counts}, the maximum magnitude of the absolute residual
  between the expected value and our calculated value is \(5.529374\);
  the maximum magnitude of the fractional residual is \(0.0028803\).
\item
  Over the \textbf{number of bins} tested and for the stacked
  \textbf{average masses}, the maximum magnitude of the absolute
  residual between the expected value and our calculated value is
  \(\ensuremath{9.57\times 10^{-4}}\); the maximum magnitude of the
  fractional residual is \(\ensuremath{6.6260105\times 10^{-5}}\).
\end{enumerate}

\begin{figure}

{\centering \includegraphics{physics_validation_files/figure-pdf/unnamed-chunk-36-1.pdf}

}

\caption{Absolute comparison of our implementation to the expected
results. Note that over-plotting may obscure the expected values, if the
approximated values are very close.}

\end{figure}

\begin{figure}

{\centering \includegraphics{physics_validation_files/figure-pdf/unnamed-chunk-37-1.pdf}

}

\caption{Relative residual implementation or our approximation to the
expected results. Points in blue pass, and points in red fail, the
absolute residual test.}

\end{figure}

\begin{figure}

{\centering \includegraphics{physics_validation_files/figure-pdf/unnamed-chunk-38-1.pdf}

}

\caption{Absolute comparison of our implementation to the expected
results for average masses. Note that over-plotting may obscure the
expected values, if the approximated values are very close.}

\end{figure}

\begin{figure}

{\centering \includegraphics{physics_validation_files/figure-pdf/unnamed-chunk-39-1.pdf}

}

\caption{Relative residual implementation or our approximation to the
expected results for average masses. Points in blue pass, and points in
red fail, the absolute residual test.}

\end{figure}

\hypertarget{cluster-mass-profile}{%
\chapter{Cluster Mass Profile}\label{cluster-mass-profile}}

The Cluster Mass profiles are calculated by a \textbf{separate CosmoSIS
module}, which is a python wrapper to Tom McClintock's Cluster\_toolkit
code. The code exists in the folder \texttt{delta\_sigma}.

\textbf{Before your run this test, make sure to go to the folder
\texttt{deltasigma} and run \texttt{cosmosis\ deltasigma\_test.ini} to
produce the proper test files.} If you are using the supplied
\texttt{GNUmakefile} to build the document, this is done automatically.

\hypertarget{xi-nfw}{%
\section{\texorpdfstring{\(\xi\)
NFW}{\textbackslash xi NFW}}\label{xi-nfw}}

\(\xi_{NFW}(r)\) is the cluster matter correlation function in the
1-halo regime, adopting an NFW model.

For this testing, we compare results from our implementation to the
function evaluations from directly calling Tom McClintock's python/c
package, for a cluster of
\(M_{200m}=3.199267137797384375e+14 M_\odot/h\) and \(c=5\).

\begin{enumerate}
\def\labelenumi{\alph{enumi}.}
\tightlist
\item
  The comparisons are made at 1000 unique values of \(r\).
\item
  Over the range tested (\(0.01 < r < 1000\)), the maximum magnitude of
  the absolute residual between the expected value and our calculated
  value is \(\ensuremath{4.7040741\times 10^{4}}\); the maximum
  magnitude of the fractional residual is \(6.6337697\).
\item
  For this test, we require the relative differences to be less than
  \(0.001\), or the absolute difference to be less than
  \(\ensuremath{10^{-12}}\).
\item
  \textbf{For this test, 297 of 1000 points pass.}
\end{enumerate}

The values of \(\xi\) cover a very wide range, so a logarithmic
\(y\)-axis is needed. However, many points in the plot have \(\xi = 0\).
In order to allow the use of a logarithmic axis, these points are
suppressed.

\begin{figure}

{\centering \includegraphics{physics_validation_files/figure-pdf/unnamed-chunk-42-1.pdf}

}

\caption{Absolute comparison of our implementation to the expected
results. Note that over-plotting may obscure the expected values, if the
approximated values are very close. Note that we suppress negative
values of \(\xi\) in this plot, so that we can use a log scale.}

\end{figure}

\begin{verbatim}
Warning: Removed 559 rows containing missing values (`geom_point()`).
\end{verbatim}

\begin{figure}

{\centering \includegraphics{physics_validation_files/figure-pdf/unnamed-chunk-43-1.pdf}

}

\caption{Relative residual implementation or our approximation to the
expected results. Points in blue pass, and points in red fail, the
absolute residual test.}

\end{figure}

\hypertarget{xi-2-halo}{%
\section{\texorpdfstring{\(\xi\)
2-halo}{\textbackslash xi 2-halo}}\label{xi-2-halo}}

\(\xi_{mm}(r)\) is the matter correlation function in the cluster 2-halo
regime, adopting a non-linear matter power spectrum, and not accounting
for cluster bias.

For this testing, we compare results from our implementation to the
function evaluations from directly calling Tom McClintock's python/c
package, for \(z=0.201010\).

\begin{enumerate}
\def\labelenumi{\alph{enumi}.}
\tightlist
\item
  The comparisons are done at 819 unique values of \(r\)
\item
  Over the range tested (\(0.01 < r < 124.192135\)), the maximum
  magnitude of the absolute residual between the expected value and our
  calculated value is \(0.1660698\); the maximum magnitude of the
  fractional residual is \(\ensuremath{6.6336039\times 10^{-5}}\).
\item
  For this test, we require the relative differences to be less than
  \(0.001\), or the absolute difference to be less than
  \(\ensuremath{10^{-12}}\).
\item
  \textbf{For this test, 819 of 819 points pass.}
\end{enumerate}

The values of \(\xi\) cover a very wide range, so a logarithmic
\(y\)-axis is needed. However, many points in the plot have \(\xi = 0\).
In order to allow the use of a logarithmic axis, these points are
suppressed.

\begin{figure}

{\centering \includegraphics{physics_validation_files/figure-pdf/unnamed-chunk-46-1.pdf}

}

\caption{Absolute comparison of our implementation to the expected
results. Note that over-plotting may obscure the expected values, if the
approximated values are very close. Note that we suppress negative
values of \(\xi\) in this plot, so that we can use a log scale.}

\end{figure}

\begin{figure}

{\centering \includegraphics{physics_validation_files/figure-pdf/unnamed-chunk-47-1.pdf}

}

\caption{Relative residual comparison or our implementation to the
expected results. Points in blue pass, and points in red fail, the
absolute residual test. Note that points for which both the predicted
and expected values are zero are suppressed.}

\end{figure}

\hypertarget{sigma-nfw}{%
\section{\texorpdfstring{\(\Sigma\)
NFW}{\textbackslash Sigma NFW}}\label{sigma-nfw}}

\(\Sigma_{NFW}(R)\) is the cluster surface mass density profile in the
1-halo regime, adopting an NFW model.

For this testing, we compare results from our implementation to the
function evaluations from directly calling Tom McClintock's python/c
package, for a cluster of
\(M_{200m}=3.199267137797384375e+14 M_\odot/h\) and \(c=5\).

\begin{enumerate}
\def\labelenumi{\alph{enumi}.}
\tightlist
\item
  The comparisons are done at 1000 unique values of \(R\).
\item
  Over the range tested (\(0.01 < R < 50.118723\)), the maximum
  magnitude of the absolute residual between the expected value and our
  calculated value is \(214.9236971\); the maximum magnitude of the
  fractional residual is \(0.1395469\).
\item
  For this test, we require the relative differences to be less than
  \(0.001\), or the absolute difference to be less than
  \(\ensuremath{10^{-12}}\).
\item
  \textbf{For this test, 3 of 1000 points pass.}
\end{enumerate}

\begin{figure}

{\centering \includegraphics{physics_validation_files/figure-pdf/unnamed-chunk-50-1.pdf}

}

\caption{Absolute comparison of our implementation to the expected
results. Note that over-plotting may obscure the expected values, if the
approximated values are very close.}

\end{figure}

\begin{verbatim}
Warning: Removed 993 rows containing missing values (`geom_point()`).
\end{verbatim}

\begin{figure}

{\centering \includegraphics{physics_validation_files/figure-pdf/unnamed-chunk-51-1.pdf}

}

\caption{Relative residual comparison or our approximation to the
expected results. Points in blue pass, and points in red fail, the
absolute residual test.}

\end{figure}

\hypertarget{sigma-2-halo}{%
\section{\texorpdfstring{\(\Sigma\)
2-halo}{\textbackslash Sigma 2-halo}}\label{sigma-2-halo}}

\(\Sigma_{mm}(R)\) is the cluster surface mass density profile in the
cluster 2-halo regime, adopting a non-linear matter power spectrum, and
not accounting for cluster bias.

\begin{enumerate}
\def\labelenumi{\alph{enumi}.}
\tightlist
\item
  The comparisons are done at 939 unique values of \(R\).
\item
  Over the range tested (\(0.01 < R < 29.790245\)), the maximum
  magnitude of the absolute residual between the expected value and our
  calculated value is \(0.0052628\); the maximum magnitude of the
  fractional residual is \(0.0014654\).
\item
  For this test, we require the relative differences to be less than
  \(0.01\), or the absolute difference to be less than
  \(\ensuremath{10^{-12}}\).
\item
  \textbf{For this test, 939 of 939 points pass.}
\end{enumerate}

\begin{figure}

{\centering \includegraphics{physics_validation_files/figure-pdf/unnamed-chunk-54-1.pdf}

}

\caption{Absolute comparison of our implementation to the expected
results. Note that over-plotting may obscure the expected values, if the
approximated values are very close.}

\end{figure}

\begin{figure}

{\centering \includegraphics{physics_validation_files/figure-pdf/unnamed-chunk-55-1.pdf}

}

\caption{Relative residual comparison or our approximation to the
expected results. Points in blue pass, and points in red fail, the
absolute residual test.}

\end{figure}

\hypertarget{halo-bias-bm_200m-z}{%
\section{\texorpdfstring{Halo bias
\(b(M_{200m}, z)\)}{Halo bias b(M\_\{200m\}, z)}}\label{halo-bias-bm_200m-z}}

\(b(M_{200m}, z)\) is the halo bias at mass \(M_{200m}\) and redshift
\(z\).

For this testing, we compare results from our implementation to the
function evaluations from directly calling Tom McClintock's python/c
package.

\begin{enumerate}
\def\labelenumi{\alph{enumi}.}
\tightlist
\item
  The comparisons are done at 2 different \(z\) values, and 20 different
  \(M_{200m}\) values.
\item
  Over the range tested (\(12 < Log M_{200m} < 15.8\)), the maximum
  magnitude of the absolute residual between the expected value and our
  calculated value is \(\ensuremath{2.32524\times 10^{-6}}\); the
  maximum magnitude of the fractional residual is
  \(\ensuremath{1.1546185\times 10^{-7}}\).
\item
  For this test, we require the relative differences to be less than
  \(0.001\), or the absolute difference to be less than
  \(\ensuremath{10^{-12}}\).
\item
  \textbf{For this test, 40 of 40 points pass.}
\end{enumerate}

Each panel in the figure below corresponds to a \(z\) bin; the value of
\(z\) is specified in the panel top label.

\begin{figure}

{\centering \includegraphics{physics_validation_files/figure-pdf/unnamed-chunk-58-1.pdf}

}

\caption{Absolute implementation of our approximation to the expected
results. Note that over-plotting may obscure the expected values, if the
approximated values are very close.}

\end{figure}

\begin{verbatim}
Scale for y is already present.
Adding another scale for y, which will replace the existing scale.
\end{verbatim}

\begin{figure}

{\centering \includegraphics{physics_validation_files/figure-pdf/unnamed-chunk-59-1.pdf}

}

\caption{Relative residual implementation of our approximation to the
expected results. Blue points pass, and red points fail, the absolute
residual test.}

\end{figure}

\hypertarget{mass-profile-corrections}{%
\chapter{Mass Profile corrections}\label{mass-profile-corrections}}

\hypertarget{mis-centering-offset}{%
\section{Mis-centering offset}\label{mis-centering-offset}}

\(P(R_{mis})\) is the probabilistic distribution of the mis-centering
offset \(R_{mis}\). Our calculation is implemented in the file
\texttt{roffset\_t.hh}.

For this testing, we compare our function output to the evaluations from
Yuanyuan Zhang in the DES Y1 centering analysis (Zhang et al.~2019).

\begin{enumerate}
\def\labelenumi{\alph{enumi}.}
\tightlist
\item
  The comparisons are done at 20 unique \(R_{mis}\) values.
\item
  Over the range tested \((0.01 < R_{mis} < 2)\), the maximum magnitude
  of the absolute residual between the expected value and our
  approximation is \(\ensuremath{8.8817842\times 10^{-16}}\); the
  maximum magnitude of the fractional residual is
  \(\ensuremath{3.9331606\times 10^{-16}}\).
\item
  For this test, we require the relative residuals to be less than
  \(\ensuremath{10^{-6}}\), or the absolute residuals to be less than
  \(\ensuremath{10^{-12}}\).
\item
  \textbf{For this test, 20 of 20 points pass.}
\end{enumerate}

\begin{figure}

{\centering \includegraphics{physics_validation_files/figure-pdf/unnamed-chunk-62-1.pdf}

}

\caption{Absolute comparison of our implementation to the expected
results. Note that over-plotting may obscure the expected values, if the
approximated values are very close.}

\end{figure}

\begin{figure}

{\centering \includegraphics{physics_validation_files/figure-pdf/unnamed-chunk-63-1.pdf}

}

\caption{Relative residual comparison or our implementation to the
expected results. Points in blue pass, and points in red fail, the
absolute residual test.}

\end{figure}

\hypertarget{lensing}{%
\chapter{Lensing}\label{lensing}}

\hypertarget{critical-density}{%
\section{Critical density}\label{critical-density}}

\(\Sigma_{crit}(z_l, z_s)\) is the inverse of the critical surface
density of a lens system, which changes with the redshift of the lens
\(z_l\) and the redshift of the source galaxy \(z_s\). Our calculation
is implemented in the file \texttt{sigma\_crit\_inverse\_t.hh}.

For this testing, we compare our function output to the evaluations
using the astropy package with a flat cosmology \(\Omega_m=0.3\),
\(\Omega_b=0.06\), \(h_0=0.7\).

\begin{enumerate}
\def\labelenumi{\alph{enumi}.}
\tightlist
\item
  The comparisons are done at, 100 unique \(z_s\) values, with 10
  different \(z_l\) values
\item
  Over the range tested (\(0.01 < zs < 1\)), the maximum magnitude of
  the absolute residual between the expected value and our calculated
  value is \(\ensuremath{4.7767481\times 10^{-17}}\); the maximum
  magnitude of the fractional residual is \(0.1347652\).
\item
  For this test, we require the relative differences to be less than
  \(\ensuremath{10^{-4}}\), \textbf{or} the absolute difference to be
  less than \(\ensuremath{10^{-21}}\).
\item
  \textbf{For this test, 460 of 1000 points pass.}
\end{enumerate}

Each panel in the figure below corresponds to a \(\lambda_{c}\) bin; the
value of the \(\lambda_{c}\) mass is specified in the panel top label.

\begin{figure}

{\centering \includegraphics{physics_validation_files/figure-pdf/unnamed-chunk-66-1.pdf}

}

\caption{Absolute comparison of our implementation to the expected
results. Note that over-plotting may obscure the expected values, if the
approximated values are very close.}

\end{figure}

\begin{figure}

{\centering \includegraphics{physics_validation_files/figure-pdf/unnamed-chunk-67-1.pdf}

}

\caption{Absolute residual comparison of our implementation to the
expected results. Blue points pass, and red points fail, the relative
residual test.}

\end{figure}

\begin{verbatim}
Warning: Removed 540 rows containing missing values (`geom_point()`).
\end{verbatim}

\begin{figure}

{\centering \includegraphics{physics_validation_files/figure-pdf/unnamed-chunk-68-1.pdf}

}

\caption{Relative residual comparison or our implementation to the
expected results. Points in blue pass, and points in red fail, the
absolute residual test. In this plot, the relative residual is forced to
be 0 when both the true and test values are 0, as is the case when zs
\textless= zl}

\end{figure}

\hypertarget{session-information}{%
\chapter{Session information}\label{session-information}}

The following R package versions were used to prepare this document.

\begin{verbatim}
R version 4.1.3 (2022-03-10)
Platform: x86_64-conda-linux-gnu (64-bit)
Running under: SUSE Linux Enterprise Server 15 SP4

Matrix products: default
BLAS/LAPACK: /global/common/software/des/common/Conda_Envs/cosmosis-global/lib/libopenblasp-r0.3.20.so

locale:
 [1] LC_CTYPE=en_US.UTF-8       LC_NUMERIC=C              
 [3] LC_TIME=en_US.UTF-8        LC_COLLATE=en_US.UTF-8    
 [5] LC_MONETARY=en_US.UTF-8    LC_MESSAGES=en_US.UTF-8   
 [7] LC_PAPER=en_US.UTF-8       LC_NAME=C                 
 [9] LC_ADDRESS=C               LC_TELEPHONE=C            
[11] LC_MEASUREMENT=en_US.UTF-8 LC_IDENTIFICATION=C       

attached base packages:
[1] stats     graphics  grDevices utils     datasets  methods   base     

other attached packages:
 [1] lubridate_1.9.2 forcats_1.0.0   stringr_1.5.0   dplyr_1.1.2    
 [5] purrr_1.0.1     readr_2.1.4     tidyr_1.3.0     tibble_3.2.1   
 [9] ggplot2_3.4.2   tidyverse_2.0.0

loaded via a namespace (and not attached):
 [1] pillar_1.9.0     compiler_4.1.3   tools_4.1.3      bit_4.0.5       
 [5] digest_0.6.31    timechange_0.2.0 jsonlite_1.8.4   evaluate_0.20   
 [9] lifecycle_1.0.3  gtable_0.3.3     pkgconfig_2.0.3  rlang_1.1.0     
[13] cli_3.6.1        parallel_4.1.3   yaml_2.3.7       xfun_0.39       
[17] fastmap_1.1.1    withr_2.5.0      knitr_1.42       generics_0.1.3  
[21] vctrs_0.6.2      hms_1.1.3        bit64_4.0.5      grid_4.1.3      
[25] tidyselect_1.2.0 glue_1.6.2       R6_2.5.1         fansi_1.0.4     
[29] vroom_1.6.1      rmarkdown_2.21   farver_2.1.1     tzdb_0.3.0      
[33] magrittr_2.0.3   scales_1.2.1     htmltools_0.5.5  colorspace_2.1-0
[37] labeling_0.4.2   utf8_1.2.3       stringi_1.7.12   munsell_0.5.0   
[41] crayon_1.5.2    
\end{verbatim}



\end{document}
